% year: תשע"ח
% type: מבחן
% semester: א
% moed: ב
% number: 2
% points: 20
% categories: continuity, func-limits, lhopital
% source: exams/מבחנים/תשעח/מועד ב תשעח.pdf

\begin{question}
מצאו את נקודות אי-הרציפות של הפונקציה
\[ f(x)=\frac{x\ln|x|}{x^2-1} \]
וקבעו את סוגיהן.
\end{question}

\begin{hint}
בנקודות $x=\pm1$ גם המונה מתאפס — השתמשו בגבול $\lim_{t\to0}\frac{\ln(1+t)}{t}=1$ או בלופיטל. ב-$x=0$ השתמשו בכך ש-$x\ln|x|\to0$.
\end{hint}

\begin{solution}
\textbf{תחום ההגדרה:} $\ln|x|$ מוגדר עבור $x\neq0$, והמכנה מתאפס עבור $x=\pm1$. לכן $f$ מוגדרת עבור $x\notin\{0,1,-1\}$, ושם היא רציפה כמנה של פונקציות רציפות עם מכנה שונה מאפס. נקודות אי-הרציפות הן $x=0,\,1,\,-1$.

\textbf{$x=0$:} ידוע ש-$\lim_{x\to0}x\ln|x|=0$ (למשל בלופיטל: $\lim_{x\to0^+}\frac{\ln x}{1/x}=\lim_{x\to0^+}\frac{1/x}{-1/x^2}=\lim_{x\to0^+}(-x)=0$, ומשמאל זה אותו דבר כי $x\ln|x|$ אי-זוגית). המכנה שואף ל-$-1$. לכן
\[ \lim_{x\to0}f(x)=\frac{0}{-1}=0. \]
הגבול קיים וסופי אך $f$ אינה מוגדרת ב-$0$: \textbf{אי-רציפות סליקה}.

\textbf{$x=1$:} בסביבת $1$, $|x|=x$, והביטוי מהצורה $\frac00$. נכתוב
\[ f(x)=\frac{x}{x+1}\cdot\frac{\ln x}{x-1}. \]
נציב $t=x-1\to0$: $\frac{\ln x}{x-1}=\frac{\ln(1+t)}{t}\to1$ (גבול מפורסם). לכן
\[ \lim_{x\to1}f(x)=\frac{1}{2}\cdot1=\frac12. \]
\textbf{אי-רציפות סליקה}.

\textbf{$x=-1$:} נשים לב ש-$f(-x)=\frac{-x\ln|x|}{x^2-1}=-f(x)$, כלומר $f$ אי-זוגית. לכן
\[ \lim_{x\to-1}f(x)=-\lim_{x\to1}f(x)=-\frac12. \]
(או ישירות: בסביבת $-1$, $f(x)=\frac{x}{x-1}\cdot\frac{\ln(-x)}{x+1}$, ועם $t=-x-1\to0$: $\frac{\ln(-x)}{x+1}=\frac{\ln(1+t)}{-t}\to-1$, ו-$\frac{x}{x-1}\to\frac12$.)
\textbf{אי-רציפות סליקה}.

\textbf{סיכום:} שלוש נקודות אי-הרציפות $x=0,1,-1$ כולן סליקות; ניתן להגדיר $f(0)=0$, $f(1)=\frac12$, $f(-1)=-\frac12$ ולקבל פונקציה רציפה על כל $\R$.
\end{solution}
